
% ---------------------------------------------------------------------------
% Abstract and keywords
% ---------------------------------------------------------------------------

\chapter*{Abstract}

ISOGEO's Scan product explores various geospatial data sources to identify
their contents and extract their metadata. This information supports the data
cataloguing process, which lies at the core of ISOGEO's activities by
facilitating the inventory, documentation and enhancement of geospatial data.
Scan supports databases, enterprise geodatabases and several geospatial file
formats.

Historically, all of these processes relied on FME workbenches, resulting in an
external dependency and additional licensing costs. ISOGEO therefore began
progressively replacing them with Python scripts integrated directly into the
product. My work-study assignment consisted in continuing this migration, which
had already been initiated for sources such as PostGIS and Esri enterprise
geodatabases.

After analysing and documenting the existing developments, I contributed to
the implementation of processing components for several sources, including
vector and raster data, GeoPackage, CAD data, Oracle, SQL Server and
GeoParquet. The most recent work also focused on point cloud data in LAS and LAZ
formats. These scripts extract the tables, fields, coordinate reference systems,
geometries, spatial extents and other metadata required by Scan. They also
generate data signatures.

GDAL/OGR was the main library used, together with additional libraries adapted
to specific formats, such as LibreDWG or even SQLAlchemy. The scripts were then
organised within a modular architecture and prepared for distribution as a
standalone Windows executable using PyInstaller. This executable is built in a
Docker environment to ensure a reproducible process independent of the libraries
installed on the developer's computer. Regression tests and signature integrity
tests are also used to validate the processing components.

This work contributes to reducing Scan's dependency on FME, facilitating the
maintenance of its processing components and strengthening their integration
into ISOGEO's software environment.

\vspace{0.7cm}

\noindent
\textbf{Keywords:}
geomatics; Scan; ISOGEO; Python; FME migration; metadata; data signature;
GDAL/OGR; GeoPackage; GeoParquet; CAD; spatial databases; point clouds; Docker;
PyInstaller.
